\color{unidarkgreen}\section[Техническое обслуживание]{Техническое обслуживание}
\color{black}

\color{unidarkgreen}\subsection{Общие указания}\color{black}

\begin{enumerate}[label=\arabic{section}.\arabic{subsection}.\arabic*, align=left, labelsep=4pt, leftmargin=0pt, itemindent=2cm]

\item
Техническое обслуживание (ТО) шкафа должно производится в соответствие с «Правилами технического обслуживания устройств и комплексов релейной защиты и автоматики» утвержденные приказом Минэнерго России от 13 июля 2020 г. № 555.

\item
Техническое обслуживание устройства РЗА представляет собой деятельность по предотвращению отказов функционирования устройства РЗА, осуществляемая при выполнении работ по настройке параметров (уставок) срабатывания (возврата), алгоритмов функционирования, периодической проверке работоспособности, выявлению причин отказов и устранению обнаруженных неисправностей устройства РЗА.

\item
Под циклом технического обслуживания устройств РЗА понимается интервал времени между двумя ближайшими техническими контролями или профилактическими восстановлениями.

\setlength{\extrarowheight}{0.06cm}
\setlength{\tabcolsep}{1.4pt}
\setlength{\LTleft}{0pt}
\setlength{\LTright}{0pt}
\footnotesize
\begin{longtable}{|>{\centering\arraybackslash}m{1.7cm}|>{\centering\arraybackslash}m{3cm}|>{\centering\arraybackslash}m{2cm}|*{26}{>{\centering\arraybackslash}m{0.3cm}|}}
\caption{Периодичность технического обслуживания устройств релейной защиты и автоматики\hfill\vspace{-0.5\baselineskip}}\label{sec3:tbl1}\\ 
\hline
\rowcolor{gray!30}
Категория помещений & Вид аппаратного исполнения устройства РЗА & Цикл технического обслуживания, лет & 
\multicolumn{26}{c|}{Количество лет эксплуатации} \\
\hline
\rowcolor{gray!30}
& & & 0 & 1 & 2 & 3 & 4 & 5 & 6 & 7 & 8 & 9 & 10 & 11 & 12 & 13 & 14 & 15 & 16 & 17 & 18 & 19 & 20 & 21 & 22 & 23 & 24 & 25 \\
\hline
\endfirsthead

\caption*{\hspace{3pt}\emph{Продолжение таблицы \ref{sec3:tbl1}\hfill\vspace{-0.5\baselineskip}}} \\
\hline
\rowcolor{gray!30}
Категория помещений & Вид аппаратного исполнения устройства РЗА & Цикл технического обслуживания, лет & 
\multicolumn{26}{c|}{Количество лет эксплуатации} \\
\hline
\rowcolor{gray!30}
& & & 0 & 1 & 2 & 3 & 4 & 5 & 6 & 7 & 8 & 9 & 10 & 11 & 12 & 13 & 14 & 15 & 16 & 17 & 18 & 19 & 20 & 21 & 22 & 23 & 24 & 25 \\
\hline
\endhead

\hline
\endfoot
\endlastfoot

% Тело таблицы:
I  & электромеханическое & 8 	& Н & К1 &   &    & К  &   &    &   & В  &    &   &   & К  &   &   &    & В  &   &    &   & К  &    &   &   & В  &   \\ \hline
   & микроэлектронное    & 6 	& Н & К1 & Т & К  & Т  & Т & В  & Т & Т  & К  & Т & Т & В  & Т & Т & К  & Т  & Т & В  & Т & Т  & К  & Т & Т & В  & Т \\ \hline
   & микропроцессорное   & 4(ТК)& Н & К1 &   &    & ТК &   &    &   & ТК &    &   &   & ТК &   &   &    & ТК &   &    &   & ТК &    &   &   & ТК &   \\ \hline
   &                     & 8(В) & Н & К1 &   &    & К  &   &    &   & В  &    &   &   & К  &   &   &    & В  &   &    &   & К  &    &   &   & В  &   \\ \hline
II & электромеханическое & 6 	& Н & К1 &   & К  &    &   & В  &   &    & К  &   &   & В  &   &   & К  &    &   & В  &   &    & К  &   &   & В  &   \\ \hline
   & микроэлектронное    & 6 	& Н & К1 & Т & К  & Т  & Т & В  & Т & Т  & К  & Т & Т & В  & Т & Т & К  & Т  & Т & В  & Т & Т  & К  & Т & Т & В  & Т \\ \hline
   & микропроцессорное   & 3(ТК)& Н & К1 &   & ТК &    &   & ТК &   &    & ТК &   &   & ТК &   &   & ТК &    &   & ТК &   &    & ТК &   &   & ТК &   \\ \hline
   &                     & 6(В) & Н & К1 &   & К  &    &   & В  &   &    & К  &   &   & В  &   &   & К  &    &   & В  &   &    & К  &   &   & В  &   \\ \hline
\multicolumn{29}{|p{\dimexpr\linewidth-2\tabcolsep\relax}|}{%
  \raggedright
  Примечание: В -- профилактическое восстановление; К -- профилактический контроль; К1 -- первый профилактический контроль; \\
  Н -- проверка при новом включении (наладка); Т -- тестовый контроль; ТК -- технический контроль.%
} \\ \hline

\end{longtable}
\setlength{\tabcolsep}{6pt}
\normalsize
\setlength{\extrarowheight}{0.0cm}

\item
Интервалы между различными видами технического обслуживания устройств РЗА, указанные в таблице, являются максимально допустимыми и могут быть сокращены по решению собственника объекта электроэнергетики с учетом условий эксплуатации и технического состояния конкретного устройства РЗА.

\end{enumerate}

\color{unidarkgreen}\subsection{Меры безопасности}\color{black}

\begin{enumerate}[label=\arabic{section}.\arabic{subsection}.\arabic*, align=left, labelsep=4pt, leftmargin=0pt, itemindent=2cm]

\item
Конструкция шкафа пожаробезопасна в соответствии с ГОСТ 12.1.004-91 и обеспечивает безопасность обслуживания в соответствии с ГОСТ IEC 61439-1-2024, ГОСТ 12.2.007.0-75.

\item
По способу защиты человека от поражения электрическим током шкаф и терминал соответствуют классу 0I по ГОСТ 12.2.007.0-75.

\item
При эксплуатации и техническом обслуживании шкафа необходимо руководствоваться действующими требованиями по охране труда (правила безопасности) и НТД.

\item
При соблюдении требований эксплуатации и хранения шкаф не создает опасность для окружающей среды.

\end{enumerate}

\color{unidarkgreen}\subsection{Порядок технического обслуживания}\color{black}

\begin{enumerate}[label=\arabic{section}.\arabic{subsection}.\arabic*, align=left, labelsep=4pt, leftmargin=0pt, itemindent=2cm]

\item
Виды работ и их последовательность должна соответствовать приложению \ref{app:pto} {\color{red}«Правила технического обслуживания устройств и комплексов релейной защиты и автоматики»}.

\item
При проверке работоспособности ИЭУ необходимо руководствоваться методиками проверок и руководством по эксплуатации на ИЭУ (при наличии - руководством на ПО для мониторинга и настройки).

\item
Обслуживание шкафа должны осуществлять специалисты, имеющие соответствующие допуски к работам в электроустановках напряжением до 1000 В и необходимую квалификацию.

\item
Для технического обслуживания шкафа должны использоваться ИО, аттестованные и поверенные в установленном порядке, внесённые в Государственный Реестр СИ.

\end{enumerate}